Neural network inference is increasingly deployed inside cloud-native software systems that must balance latency, deployability, scalability, and isolation. Splitting a model across microservices can improve modularity and make security boundaries more explicit, but it also forces intermediate activations to cross service interfaces. Each boundary adds serialization, transport, scheduling, and coordination cost. The central problem studied in this thesis is therefore not only where to split a neural network, but whether a decomposition that appears acceptable in a controlled local setting remains credible once it is moved into Kubernetes, Azure, and security-hardened deployment contexts.

This thesis addresses that problem through a staged experimental evaluation of ResNet-18 inference. A proof-of-concept gRPC-based microservice framework is used to compare monolithic inference, two-part model partitions, and synchronously chained multi-service deployments. The evaluation first screens coarse architectural split points under controlled local execution, then refines the selected region and explains the observed behaviour using activation-transfer size and compute distribution. The same decomposition family is then evaluated in local Kubernetes and Azure Kubernetes Service, after which the selected Azure deployment is hardened with service identity and mutual TLS. Confidential-execution mechanisms are treated as a later incremental protection layer on top of the secured Azure path.

The results show that monolithic execution remains the fastest baseline, but that not all microservice decompositions are equally costly. Early splits suffer from large activation transfers, while very late splits can become compute-degenerate. Mid-network boundaries provide the most credible two-part trade-off. Fine-grained refinement within this region does not materially improve latency, supporting the use of coarse architectural boundaries for later stages. In Kubernetes and Azure, latency increases with service count, but the increase is bounded and non-linear: later boundaries add less cost when the transferred activations are smaller. Azure preserves the condition ordering observed in local Kubernetes, indicating directional transfer despite different absolute latencies. Adding mutual TLS and service identity introduces measurable but bounded overhead, alongside clear operational complexity from sidecars, policies, and readiness delays.

Overall, the thesis concludes that microservice-based inference is viable only when architectural partitioning, deployment context, and security techniques are evaluated together. The main engineering trade-off is not between monolithic and distributed inference in isolation, but between deployability and stronger isolation on one side and cumulative boundary-related overhead on the other.
